\section{RESULTS}
\label{sec:results}

This section reports the three measurements set up in
Section~\ref{subsec:method-eval}. Parser accuracy on OmniDocBench tests M1;
translation quality and length-compliance on WMT24++ test M2 and M3. Full
per-slice tables are kept in Appendices~\ref{app:parser-slices}
and~\ref{app:translation-full}.

\subsection{Finding 1: The Parser Loses Little, but What It Loses Is
Unrecoverable}
\label{subsec:res-parser}

\emph{The parser usually recovers page regions, labels and order. Its largest
apparent localisation weakness is over-segmentation, which is mostly benign for
translation; the dangerous errors are missed regions, weak OCR slices and missed
structured content.}

Table~\ref{tab:res-headline} gives the global scores over all 1{,}651
OmniDocBench pages using the split parser configuration and fine ground-truth
granularity. The granularity-invariant matcher reaches $F_1@.5=0.940$, showing
that predicted and annotated regions usually cover the same page areas. The
strict COCO 1-1 score is lower ($0.654$), mainly because sparse-block splitting
emits multiple useful fragments for a single annotated block. That gap is the
cost of granularity, not simply missed content.

\begin{table}[H]
\centering
\caption{Headline structure-analysis results over all 1{,}651 pages.}
\label{tab:res-headline}
\small
\begin{tabular}{@{}l l r@{}}
\toprule
\textbf{Task} & \textbf{Metric} & \textbf{Value} \\
\midrule
Localisation      & $F_1@.5$ ($N$--$M$)          & 0.940 \\
                  & $F_1@.5$ (COCO 1-1)         & 0.654 \\
                  & mF1$[.5{:}.95]$ ($N$--$M$)   & 0.700 \\
Classification    & label accuracy              & 0.889 \\
Text recognition  & edit distance ($\downarrow$) & 0.105 \\
                  & CER ($\downarrow$)          & 0.108 \\
Reading order     & NED ($\downarrow$)          & 0.067 \\
Formula           & coverage                    & 0.917 \\
                  & CDM $F_1$                   & 0.863 \\
Table (content)   & coverage                    & 0.801 \\
                  & CER on matched ($\downarrow$) & 0.216 \\
\bottomrule
\end{tabular}
\end{table}

The remaining risks are concentrated rather than diffuse
(Table~\ref{tab:res-risk-slices}). English OCR is strong (CER $0.055$), while
Simplified and Traditional Chinese are substantially harder; degraded historical
documents are the clearest failure case. Formula recognition is better than
edit distance alone suggests: CDM $F_1=0.863$ indicates many string differences
render equivalently, so the main formula issue is detection coverage, especially
on three-column pages. Tables show the same pattern: matched content is readable,
but about one table in five is not covered at all.

\begin{table}[H]
\centering
\caption{Parser failure slices that most affect downstream translation.}
\label{tab:res-risk-slices}
\small
\setlength{\tabcolsep}{5pt}
\begin{tabularx}{\textwidth}{@{}l l r >{\raggedright\arraybackslash}X@{}}
\toprule
\textbf{Slice} & \textbf{Metric} & \textbf{Value} & \textbf{Interpretation} \\
\midrule
English text & CER $\downarrow$ & 0.055 & Intended source direction is comparatively safe. \\
Simplified Chinese text & CER $\downarrow$ & 0.224 & Script-specific OCR errors can propagate fluently. \\
Traditional Chinese text & CER $\downarrow$ & 0.403 & Small slice, but high recognition risk. \\
Historical documents & $N$--$M$ recall & 0.698 & Nearly a third of regions can be absent. \\
Three-column formulas & coverage & 0.444 & More than half of formulas are missed. \\
All tables & coverage & 0.801 & Table loss is mainly detection coverage. \\
\bottomrule
\end{tabularx}
\end{table}

\paragraph{Cell geometry.} OmniDocBench has no per-cell boxes, so the cell stage is scored
separately on PubTables-1M (Table~\ref{tab:res-cells}). Given a correct table region,
cells are usually found but their edges are not tight: $F_1$ falls from $0.902$ at IoU
$0.5$ to $0.750$ at $0.75$. The loss is again concentrated---recall is $0.945$ on ordinary
cells but $0.504$ on cells spanning several rows or columns---and again it is a loss of
geometry rather than of content, because text falling outside every detected cell is
recovered through a supplementary text-line path. What is lost is that a spanning header
reaches the translator as several fragments rather than one region.

\begin{table}[H]
\centering
\caption{Cell geometry on 1{,}499 PubTables-1M tables, given the annotated table region.}
\label{tab:res-cells}
\small
\begin{tabular}{@{}l r@{}}
\toprule
\textbf{Metric} & \textbf{Value} \\
\midrule
$F_1$ @ IoU $0.50$            & 0.902 \\
$F_1$ @ IoU $0.75$            & 0.750 \\
AP@$[.5{:}.95]$               & 0.529 \\
\midrule
Recall, ordinary cells        & 0.945 \\
Recall, spanning cells        & 0.504 \\
\bottomrule
\end{tabular}
\end{table}

Most of that gap is a deliberate trade-off rather than an oversight. The parser discards
any cell box that encloses a smaller one, which guarantees that cell regions never
overlap; removing that filter raises spanning-cell recall to $0.621$ but leaves nested
boxes in $43\%$ of tables, and nested boxes cannot be reconstructed without drawing the
same text twice, because each cell is masked and re-typeset independently. Layout
preservation therefore constrains structure analysis, and not only the reverse---the
coupling H2 asserts, seen from the opposite direction. Appendix~\ref{app:cells} reports
how sensitive these numbers are to the settings involved.

\subsection{Finding 2: Model Choice Is the Smaller Lever}
\label{subsec:res-mt}

\emph{Across five hosted LLMs, the controllable model-choice spread is narrow
compared with the spread across target languages.}

Over all 55 \texttt{en}$\rightarrow$\texttt{xx} WMT24++ pairs
(Table~\ref{tab:comet-per-lang}),
\texttt{gemini-3.1-flash-lite} has the highest mean COMET-DA ($0.841$),
followed by \texttt{gpt-5.4-nano} ($0.826$),
\texttt{gemini-2.5-flash-lite} ($0.821$), \texttt{gpt-4.1-nano} ($0.800$) and
\texttt{qwen3.5-flash} ($0.788$). The whole system band is therefore $0.053$
COMET-DA.

\begin{table}[H]
    \centering
    \small
    \caption{COMET-DA (chrF++ in parentheses) by target language $\times$
    system for representative languages. Full results are in
    Table~\ref{tab:appendix-comet-full}.}
    \label{tab:comet-per-lang}
    \setlength{\tabcolsep}{4pt}
    \makebox[\textwidth][c]{%
    \begin{tabular}{@{}lccccc@{}}
        \toprule
        Language & gemini-2.5 & gemini-3.1 & gpt-4.1 & gpt-5.4 & qwen3.5 \\
         & flash-lite & flash-lite & nano & nano & flash \\
        \midrule
        Vietnamese (vi\_VN)         & 0.844 (51.7) & \textbf{0.853 (51.4)} & 0.823 (48.5) & 0.837 (48.7) & 0.813 (45.8) \\
        German (de\_DE)             & 0.827 (56.4) & \textbf{0.829 (54.6)} & 0.811 (53.5) & 0.823 (55.2) & 0.760 (44.9) \\
        Simplified Chinese (zh\_CN) & 0.835 (33.1) & \textbf{0.852 (31.9)} & 0.828 (28.5) & 0.833 (29.0) & 0.713 (14.4) \\
        Japanese (ja\_JP)           & 0.867 (29.0) & \textbf{0.873 (27.8)} & 0.859 (27.3) & 0.868 (29.3) & 0.843 (25.8) \\
        Egyptian Arabic (ar\_EG)    & 0.769 (34.6) & \textbf{0.774 (35.8)} & 0.739 (33.1) & 0.769 (33.2) & 0.742 (28.7) \\
        Hindi (hi\_IN)              & 0.733 (38.9) & \textbf{0.751 (38.6)} & 0.703 (34.4) & 0.738 (37.4) & 0.698 (33.4) \\
        \midrule
        \textbf{Mean over 55 languages} & 0.821 (48.4) & \textbf{0.841 (48.9)} & 0.800 (44.6) & 0.826 (47.6) & 0.788 (41.7) \\
        \bottomrule
    \end{tabular}}
\end{table}

The language spread is larger: averaged over systems, the lowest-scoring
languages include Marathi ($0.693$), Hindi ($0.725$), Serbian ($0.734$), Zulu
($0.744$) and Swahili ($0.757$), while Finnish, Norwegian, Japanese, Korean and
Indonesian all exceed $0.855$. The $\approx0.18$ gap between language extremes
is more than three times the typical spread between systems on the same
language. Script alone does not explain the result, since Japanese and Korean
are non-Latin and morphologically distant from English yet remain in the top
group.

\subsection{Finding 3: A Prompt-Level Length Constraint Does Not Buy Layout
Fidelity}
\label{subsec:res-compliance}

\emph{A $\pm15\%$ character-length instruction is only weakly obeyed, mostly
measures script density, and is essentially unrelated to adequacy.}

\begin{table}[H]
    \centering
    \small
    \caption{Segment-level compliance with the $\pm15\%$ length constraint.}
    \label{tab:compliance-results}
    \setlength{\tabcolsep}{5pt}
    \begin{tabular}{@{}lccccc@{}}
        \toprule
        System & compl.\ & too & too & med.\ & IQR \\
               & \%$\uparrow$ & short & long & ratio & $\downarrow$ \\
        \midrule
        gemini-2.5-flash-lite & 71.4 & 16.4 & 12.2 & \textbf{1.000} & \textbf{0.182} \\
        gemini-3.1-flash-lite & \textbf{73.0} & 19.0 & 8.0 & 0.980 & 0.187 \\
        gpt-4.1-nano          & 68.9 & 17.9 & 13.2 & \textbf{1.000} & 0.192 \\
        gpt-5.4-nano          & 68.6 & 13.2 & 18.2 & 1.036 & 0.188 \\
        qwen3.5-flash         & 71.7 & 20.7 & 7.6 & 0.977 & 0.185 \\
        \bottomrule
    \end{tabular}
\end{table}

No model exceeds 73\% compliance (Table~\ref{tab:compliance-results}), so about
27--31\% of valid segments fall outside the requested band; the median and IQR
columns report the per-segment target/source character-length ratio and its
spread (a value near $1.0$ means target length tracks the source). Per-language results in Appendix~\ref{app:compliance}
show that the violation direction is mainly a property of target-script
character density: CJK languages are too short in 89--99\% of segments, while
German, French and Filipino violate mostly by expansion. Vietnamese expands only
mildly in characters despite being grouped a priori with German. The correlation
between per-language violation rate and COMET-DA is $r=0.066$, effectively zero.

\subsection{Where the Bottleneck Is}
\label{subsec:res-bottleneck}

The three findings are not on a single numeric scale, but they differ in kind.
Changing among the five translation systems moves adequacy within a bounded band
on content that exists. In contrast, an undetected formula, table or page region
never reaches the translator, and an OCR error becomes fluent wrong input. In
this implementation, the most plausible controllable headroom is therefore
parser coverage on hard layouts, script-specific OCR and reconstruction metrics
that are computed from the page itself rather than requested from the language
model.

\subsection{Pilot Runtime on Seven PDFs}
\label{subsec:res-pilot-runtime}

Table~\ref{tab:pipeline-pilot-runtime} profiles the complete pipeline on seven
held-out PDFs (207 pages). The set is small and unannotated, so it is an
implementation diagnostic rather than a benchmark or SLA.

Structure analysis accounts for 65.3\% of elapsed time, translation for 27.4\%
and rendering for 7.3\%. This matches the quality-side result: the visual parser
is both the largest local cost and the stage whose failures are least
recoverable.

\begin{table}[H]
\centering
\caption{Pilot full-pipeline runtime. Times are wall-clock seconds.}
\label{tab:pipeline-pilot-runtime}
\small
\setlength{\tabcolsep}{5pt}
\begin{tabular}{@{}l r r r r r r@{}}
\toprule
Document & Pages & Parse & Translate & Render & Total & s/page \\
\midrule
Book --- EN      & 27 & 50.17 & 21.41 & 4.28 & 75.87 & 2.81 \\
Book --- VI      & 36 & 36.06 & 18.41 & 3.52 & 57.98 & 1.61 \\
Mix Type --- EN  & 25 & 61.36 & 21.31 & 8.08 & 90.75 & 3.63 \\
Mix Type --- VI  & 18 & 28.20 & 20.90 & 3.11 & 52.21 & 2.90 \\
Paper --- EN     & 11 & 23.88 & 23.71 & 2.92 & 50.51 & 4.59 \\
Slide --- EN     & 32 & 68.19 & 12.79 & 8.48 & 89.46 & 2.80 \\
Slide --- VI     & 58 & 58.93 & 18.55 & 5.93 & 83.41 & 1.44 \\
\midrule
\textbf{Total / mean} & \textbf{207} & \textbf{326.79} & \textbf{137.08} & \textbf{36.32} & \textbf{500.19} & \textbf{2.42} \\
\bottomrule
\end{tabular}
\end{table}

\subsection{End-to-End Qualitative Behaviour}

Figure~\ref{fig:qualitative} shows the complete pipeline on a slide, a book page
and a mathematics-heavy page. Source geometry, figures and bypassed regions are
preserved, while translated text is re-typeset in the original regions with
group-consistent font sizes and pixel-estimated appearance. Additional examples
are in Appendix~\ref{app:examples}.

\begin{figure}[H]
  \centering
  \begin{subfigure}[t]{0.3\textwidth}
    \centering
    \includegraphics[width=\linewidth]{source_1.png}
    \caption{Slide --- source}
  \end{subfigure}\hfill
  \begin{subfigure}[t]{0.3\textwidth}
    \centering
    \includegraphics[width=\linewidth]{source_3.png}
    \caption{Book --- source}
  \end{subfigure}\hfill
  \begin{subfigure}[t]{0.3\textwidth}
    \centering
    \includegraphics[width=\linewidth]{source_math.png}
    \caption{Mathematics --- source}
  \end{subfigure}

  \vspace{0.6em}

  \begin{subfigure}[t]{0.3\textwidth}
    \centering
    \includegraphics[width=\linewidth]{des_1.png}
    \caption{Slide --- translated}
  \end{subfigure}\hfill
  \begin{subfigure}[t]{0.3\textwidth}
    \centering
    \includegraphics[width=\linewidth]{des_3.png}
    \caption{Book --- translated}
  \end{subfigure}\hfill
  \begin{subfigure}[t]{0.3\textwidth}
    \centering
    \includegraphics[width=\linewidth]{des_math.png}
    \caption{Mathematics --- translated}
  \end{subfigure}
  \caption{End-to-end output on three document genres. Top row: source pages.
  Bottom row: the corresponding pages produced by the pipeline.}
  \label{fig:qualitative}
\end{figure}
